\subsection{Theorem-Style Results and Formal Statements}

Formal environments advertise the role and authority of a block before the reader examines its details.

\begin{convention} \label{conv:theorem-style-result-hierarchy} \hlt{Theorem-Style Result Hierarchy}\\
The theorem-style category contains exactly the following proved results.
\begin{enumerate}[label=(\roman*)]
\setlength{\itemsep}{0pt}
\item \hlt{Theorem}: a major result that advances a principal objective of the surrounding work.
\item \hlt{Proposition}: an independently useful result of smaller local weight.
\item \hlt{Lemma}: a subsidiary result introduced primarily to support a later proof.
\item \hlt{Corollary}: a result requiring little additional argument after an earlier result.
\end{enumerate}
The classification communicates structural role rather than a universal rank of importance. Definitions introduce objects or terminology, while conjectures record unproved claims. This hierarchy follows standard mathematical-writing guidance \citep{mit1998mathematicalwriting}.
\end{convention}

\begin{process} \label{proc:theorem-statement-audit} \hlt{Theorem-Statement Construction and Audit}\\
Given a candidate result and its supporting argument, produce a statement that can be cited and proved without hidden context.
\begin{enumerate}[label=(\roman*)]
\setlength{\itemsep}{0pt}
\item \hlt{Purpose}: identify the reusable conclusion and its role in the dependency chain.
\item \hlt{Objects}: bind every ambient set, structure, parameter, map, and index before its first use.
\item \hlt{Hypotheses}: state the exact domain, regularity, finiteness, independence, or boundary conditions needed.
\item \hlt{Conclusion}: state existence, uniqueness, characterization, bound, convergence, or other claim explicitly.
\item \hlt{Claim census}: separate conclusions whose removal would erase a distinct proof obligation.
\item \hlt{Role}: choose theorem, proposition, lemma, or corollary by Convention \ref{conv:theorem-style-result-hierarchy}.
\item \hlt{Locality}: remove motivation, examples, proof machinery, and unused assumptions from the statement.
\item \hlt{Traceability}: attach a stable title and label, and cite external ownership of a nonoriginal result.
\end{enumerate}
The final statement should remain intelligible when read apart from the preceding paragraph. Self-contained statements and first-use symbol definitions are central expository checks \citep{knuth1988mathematicalwriting}.
\end{process}

\begin{principle} \label{princ:claim-authority-separation} \hlt{Claim--Authority Separation}\\
A formal block should display its actual epistemic status. A definition stipulates meaning; an assumption declares a condition; a conjecture remains unproved; a theorem-style result requires a proof; and an empirical statement requires data and a protocol. A citation identifies prior ownership but does not replace a proof when the present work claims to establish the result.
\end{principle}

\begin{warning} \label{warn:statement-proof-contamination} \hlt{Statement--Proof Contamination}\\
Do not place a proof plan, auxiliary construction, example, or claim of obviousness inside a theorem-style statement. Such material obscures the logical contract, creates unused hypotheses, and makes later citation unreliable. Move proof-specific objects into the proof and explanatory material after the formal owner.
\end{warning}
